\noindent\textbf{Scope and editorial standard.}
This chapter covers all 30 catalogue records, 530--559, and all 45 manuscript
pages P000981--P001025. There are 26 talk records, two headings, one plan and
one intervention recorded separately. The occurrence index contains 210
items; the problem index contains 17 definite questions, including the reused
GP-1979-41. The manuscript order is retained, including 549 before 548 and
the Fuchs intervention 551 beside the Haefliger material 553. Catalogue page
windows are preserved in the talk metadata, while summaries and item pages
identify actual boundaries. All current-status statements are assessed as of
11 November 2026. Notebook wording, verified literature and editorial
reconstruction are separated throughout.

\begin{center}
\begin{tabular}{clrr}
\toprule
Class&Historical assessment&Items&Share\\
\midrule
A&Speaker/Gelfand publication soon afterwards&31&14.8\%\\
B&Established subsequent development by others&9&4.3\%\\
C&Partial development; a precise remainder&5&2.4\%\\
D&No development identified&0&0.0\%\\
E&Known before the talk&149&71.0\%\\
F&Exact assertion or fate not identifiable&16&7.6\%\\
\bottomrule
\end{tabular}
\end{center}
\noindent Percentages are rounded. A class records the fate of the identifiable
mathematical thread, not the correctness of every notebook formula. Thirty-nine
occurrences are false as written; most require a sign, normalization, domain
or quantifier correction rather than a new research problem.

\paragraph{The year's main lines.}
The autumn concentrates on Lie algebras, quivers, multiplicities and
geometric representation theory. Serre relations, classical incidence
geometry, Gelfand--Tsetlin patterns and Verma modules are interwoven with
recent results. Abeasis--Del Fra--Kraft supply a close quiver-publication
match; the Kazhdan--Lusztig multiplicity discussion continues earlier
seminar questions. Elashvili's short classification announcement was received
only a week after his talk. The Jantzen discussion requires a particularly
important distinction: dominant deformation directions have the established
invariance property, whereas an unrestricted regular direction can change
the layer multiplicities. A complete exact computation reproduces the
published type-$A_3$ counterexample.

The winter and spring move through nonstandard perturbation theory,
integrable systems, discrete curvature, microlocal analysis, diffeomorphism
groups and foliations. M\"uller's curvature talk is matched to
Cheeger--M\"uller--Schrader, with its early announcement distinguished from
the 1984 full paper. Convergence of curvature measures does not prove
convergence of solutions of the Regge equations. McDuff's parity statements
and Milnor's discrete-Lie-group results have close subsequent publication
matches. The finite-coefficient Friedlander--Milnor comparison remains a
precise conjectural thread; broad proof claims are not treated as established
merely because a preprint exists. Kirillov and Givental lead into the final
two Wodzicki talks on index theory, eta invariants and spectral cuts.

\paragraph{What is developed here.}
Two C questions permit bounded proofs, without claims of historical priority.
For Haefliger's action of \(T(x)=x+1\) and \(S(x)=-1/x\) on
\(\mathbb {RP}^1\), every continuous function is proved to be
\(u-u\circ T+v-v\circ S\) with continuous \(u,v\). The proof gives a
uniform-norm existence estimate. It does not establish the notebook's
smooth surjectivity or smooth density; the missing distributional and
derivative estimates are stated explicitly.

For Goldsheid's coordinate-product approximation of a real Lagrangian plane,
an explicit Fourier construction disproves the proposed orthogonal-projector
gap bound \(1/\sqrt2\) already for a six-plane in \(\mathbb R^{12}\).
An odd Fourier family forces every dimension-independent constant to be at
least \(\cos(\pi/8)\). A separate argument gives the fixed-dimension upper
bound \(\cos(\pi/(2n))\), sharp for \(n=2\). The existence of any universal
constant in \([\cos(\pi/8),1)\) is not decided. Independent derivations and
the supplied computational checks corroborate these proofs; numerical
minimization is never used as a certificate of a global optimum.

\paragraph{Earlier connections and critical findings.}
The chapter links the KL thread to GP-1979-41, primitive enveloping algebras
to the 1964 occurrences, torus convexity to 1980-508.15, and eta regularity
to 1972-296.2--3 and 1972-296.5. Gilkey's no-filling regularity theorem and
Wodzicki's February 1982 publication precede the April talks; publication
dates therefore prevent a simple narrative of new seminar discoveries.
The final notes also omit a symplectic quotient in a multiplicity formula,
a gamma normalization in a fractional integral, signs in a vector-calculus
matrix, and powers of the shift parameter in a zeta identity. The correction
log separates these faithful notebook errors from doubtful readings.
Spectral-cut independence at zero does not justify independence at every
negative integer. These distinctions, and the 16 genuinely underdetermined
items, are part of the critical record rather than gaps filled by conjecture.
